\documentclass{book}

\usepackage{amssymb}
\usepackage{amsmath}
\usepackage{amsthm}
\usepackage[english]{babel}
\usepackage{tikz}
\usepackage{amsfonts}

\usepackage{hyperref}
\hypersetup{
    colorlinks=true,
    linkcolor=red,
    filecolor=magenta,      
    urlcolor=cyan
    }
\usepackage[toc]{glossaries}
\usepackage[toc,page]{appendix}

\frontmatter

\DeclareMathOperator{\cis}{cis}
\DeclareMathOperator{\cosec}{cosec}
\DeclareMathOperator{\sech}{sech}
\DeclareMathOperator{\arsinh}{arsinh}
\DeclareMathOperator{\arcosh}{arcosh}
\DeclareMathOperator{\artanh}{artanh}

\DeclareMathOperator{\Var}{Var}
\DeclareMathOperator{\poisson}{Po}
\DeclareMathOperator{\normal}{N}
\DeclareMathOperator{\binomial}{B}
\DeclareMathOperator{\real}{Re}
\DeclareMathOperator{\imag}{Im}

\DeclareMathOperator{\N}{\mathbb{N}}
\DeclareMathOperator{\Z}{\mathbb{Z}}
\DeclareMathOperator{\Q}{\mathbb{Q}}
\DeclareMathOperator{\R}{\mathbb{R}}
\DeclareMathOperator{\C}{\mathbb{C}}

\DeclareMathOperator{\bigO}{\mathcal{O}}

\title{Further Maths Revision Notes}
\author{C. Irving}
\date{}

%\begin{glossary_entries}
\makenoidxglossaries
\newglossaryentry{imaginary_constant}{
    name=imaginary constant,
    description={The square root of -1, written $i$},
    plural={imaginary constants}
}
\newglossaryentry{real_number}{
    name=real number,
    description={A number on the real number line; a number in the interval $(-\infty, \infty)$.
    The set of real numbers is written $\R$},
    plural={real numbers}
}
\newglossaryentry{complex_number}{
    name=complex number,
    description={A number in the form $a+bi$, where $a,b\in\R$.
    The set of complex numbers is written $\C$},
    plural={complex numbers}
}
\newglossaryentry{polynomial}{
    name=polynomial,
    description={An expression in the form $a_0+a_1x^1+a_2x^2+\cdots+a_nx^n$},
    plural={polynomials}
}
\newglossaryentry{complex_conjugate}{
    name=complex conjugate,
    description={A complex number with a flipped sign before its imaginary part.
    The complex conjugate of a complex number $z$ is written $z^*$},
    plural={}
}
\newglossaryentry{root}{
    name=root,
    description={A value that satisfies an equation, for example 1 and -1 are the roots of the equation $x^2-1=0$, since for $x=1$ and $x=-1$, the equation holds},
    plural={roots}
}
\newglossaryentry{complex_root}{
    name=complex root,
    description={A root $x$ of an equation where $x\in\C$.
    It is often also used to mean that $x\notin\R$},
    plural={complex roots}
}
\newglossaryentry{cartesian_graph}{
    name=Cartesian graph,
    description={A graph in the $x$-$y$ plane},
    plural={Cartesian graphs}
}
\newglossaryentry{realising_the_denominator}{
    name=realising the denominator,
    description={The process of multiplying a complex fraction in the form $\frac{w}{z}$ by $\frac{z^*}{z^*}$ (where $w,z\in\C$) to make the denominator a real value without changing the value of the fraction.},
}
\newglossaryentry{quadratic}{
    name=quadratic,
    description={A second degree \gls{polynomial}, i.e. an equation in the form $ax^2+bx+c=0$, where $a\neq0$.
    When plotted on a \gls{cartesian_graph}, it forms a parabola.
    Its \glspl{root} are given by the quadratic formula, $x=\frac{-b\pm\sqrt{b^2-4ac}}{2a}$},
    plural={quadratics}
}
\newglossaryentry{conjugate_pair}{
    name=conjugate pair,
    description={A pair of \glspl{complex_number}, where one is the \gls{complex_conjugate} of the other},
    plural={conjugate pairs}
}
\newglossaryentry{cubic}{
    name=cubic,
    description={A 3rd degree \gls{polynomial}, i.e. an equation in the form $x^3+bx^2+cx+d=0$, where $a\neq0$},
    plural={conjugate pairs}
}
\newglossaryentry{argand_diagram}{
    name=Argand diagram,
    description={A 2D graph of the complex plane},
    plural={Argand diagrams}
}
\newglossaryentry{real_part}{
    name=real part,
    description={The real term of a \gls{complex_number}, i.e. where $z=a+bi,\ a,b\in\Relbar$, the real part of $z$ is $a$
    The real part of a \gls{complex_number} $z$ is written as $real(z)$ or $\Re(z)$},
    plural={}
}
\newglossaryentry{imaginary_part}{
    name=imaginary part,
    description={The coefficient of the imaginary term of a complex number, i.e. where $z=a+bi,\ a,b\in\R$, the imaginary part of $z$ is $b$.
    The imaginary part of a \gls{complex_number} $z$ is written as $imag(z)$ or $\Im(z)$},
    plural={}
}
\newglossaryentry{modulus_argument_form}{
    name=modulus-arguemnt form,
    description={The form $r(\cos\theta+i\sin\theta)\equiv\cis\theta$},
    plural={}
}
\newglossaryentry{radian}{
    name=radian,
    description={An angle unit, equal to the angle required to make the arc length of a unit circle equal 1.
    A whole circle is $2\pi$ radians; half a circle is $\pi$ radians.
    The symbol for radians is either $^\text{c}$ or $\text{rad}$, as in $\pi^\text{c}$ or $\pi\ \text{rad}$, or is often omitted},
    plural={radians}
}
\newglossaryentry{modulus}{
    name=modulus,
    description={The magnitude of a number, where $|x|$ is the modulus of $x$.
    For positive \glspl{real_number}, the modulus is itself, for negative \glspl{real_number}, the modulus is the positive version of itself ($|x|=-x,\ x\in\R^-$).
    For a \gls{complex_number} $z$, $|z|=\sqrt{\real^2(z)+\imag^2(z)}$}, % TODO: come back
    plural={}
}
\newglossaryentry{quadrant}{
    name=quadrant,
    description={A quarter of a graph; the 1st quadrant is the top right, the second is the top left, the third is the bottom left, and the fourth is the bottom right.},
    plural={quadrants}
}
\newglossaryentry{argument}{
    name=argument,
    description={The angle that a \gls{complex_number} makes with the positive real number line.
    It is written as $\arg(z)$.
    The argument is in the real interval $(-\pi,\pi]$},
    plural={arguments}
}
\newglossaryentry{}{
    name=,
    description={},
    plural={}
}
%\end{glossary_entries}

\newtheorem{theorem}{Theorem}[section]
\newtheorem{corollary}{Corollary}[theorem]
\newtheorem{lemma}[theorem]{Lemma}
 
\begin{document}
\maketitle

\tableofcontents

\mainmatter

\chapter{Core Pure Year I} \label{ch:core_pure_i}
    \section{Complex Numbers I} \label{complex_numbers_i}
        \subsection{The Imaginary Constant}
            Using only the \glspl{real_number}, there is no number $x$ that satisfies the equation $x^2+1=0$.
            To satisfy this \gls{polynomial}, you must first define a square root of -1.
            This constant is called the \gls{imaginary_constant} $i$, defined such that $i^2=-1$.
            
            Using the \gls{imaginary_constant}, \glspl{complex_number} can be defined as $a+bi$, where $a,b\in\R$.
            The set of all numbers in the form $a+bi$ is called the set of \glspl{complex_number}, written $\C$.

            The functions $\real(z)$ and $\imag(z)$ (sometimes written as $\Re(z)$ and $\Im(z)$) take the \gls{real_part} and \gls{imaginary_part} of a \gls{complex_number} respectively.
            For example,
            \begin{equation} \label{eq:real_part}
                \text{Re}(a+bi)\equiv\Re(z):=a,\ a,b\in\R
            \end{equation}
            \begin{equation} \label{eq:imaginary_part}
                \text{Im}(a+bi)\equiv\Im(z):=bi,\ a,b\in\R
            \end{equation}

        \subsection{Complex Arithmetic}
            \Glspl{complex_number} can be added and multiplied similarly to how binomials would be.
            Given two \glspl{complex_number} $w:=a+bi$ and $z:=c+di$, where $a,b,c,d\in\R$, the following rules define complex addition and multiplication:
            \begin{equation} \label{eqn:complex_addition}
                w+z=(a+bi)+(c+di)=(a+c)+(b+d)i
            \end{equation}
            \begin{equation} \label{eqn:complex_multiplication}
                w\times z=(a+bi)(c+di)=(ac-bd)+(ad+bc)i
            \end{equation}

        \subsection{Complex Conjugates}
            The \gls{complex_conjugate} of a \gls{complex_number} of $z$, written as $z^*$ is $z$ with a flipped sign before the imaginary part.
            \begin{equation}
                (a+bi)^*:=a-bi,\ a,b\in\R
            \end{equation}
            From this definition, we can see that the sum and product of a \gls{complex_number} and its conjugate are both real.
            Let $z:=a+bi$, where $a,b\in\R$.
            \begin{equation} \label{eqn:sum_of_conjugate}
                z+z^*=2a\in\R
            \end{equation}
            \begin{equation} \label{eqn:product_of_conjugate}
                zz^*=a^2+b^2\in\R
            \end{equation}
            Using the multiplication rule in equation \ref{eqn:product_of_conjugate}, a fraction of \glspl{complex_number} can be `realised' by multiplying by the \gls{complex_conjugate} of the denominator over itself.
            \Gls{realising_the_denominator} means rewriting a fraction of \glspl{complex_number} as a fraction with a real denominator.
            \begin{equation} \label{eqn:realise_denominator}
            \begin{split}
                \frac{a+bi}{c+di}&=\frac{a+bi}{c+di}\times\frac{c-di}{c-di}\\
                &=\frac{(a+bi)(c+di)}{c^2+d^2}\\
                &=\frac{(ac-bd)+(ad+bc)i}{c^2+d^2}
            \end{split}
            \end{equation}
        \subsection{Complex Roots of Polynomials}
            All $n$th degree \glspl{polynomial} (\glspl{polynomial} with a power of $x$ up to $x^n$) have exactly $n$ \glspl{root}, counting \glspl{complex_root} and algebraic multiplicities.
            Therefore, if an equation does not cross $y=0$ on a \gls{cartesian_graph}, it must have other, complex \glspl{root}.
            The simplest example of this is the \gls{quadratic} $x^2+1=0$.
            \[x^2=-1\]
            \[x=\pm i\]
            For \glspl{quadratic}, whenever the discriminant ($\Delta:=b^2-4ac$ for a \gls{quadratic} $ax^2+bx+c$) is less than $0$, the \gls{quadratic} has \glspl{complex_root}, since the quadratic formula (\ref{eq:quadratic_formula}) can be rewritten as
            \[x=\frac{-b\pm\sqrt{\Delta}}{2a}\]
            which would give a \gls{complex_number} whenever $\Delta<0$, since $\sqrt{x}$ is complex iff $x<0$ ($x\in\R$).
            
            If a \gls{polynomial} has a \gls{complex_root} $\alpha$, it must also have the \gls{complex_conjugate} \gls{root} $\alpha^*$.
            For example, if a \gls{quadratic} has a \gls{root} $\alpha=2+3i$, then the other \gls{root} must be $\alpha^*=2-3i$.
            For a \gls{cubic}, there can only be one pair of \glspl{root}, thus the third \gls{root} cannot be a \gls{complex_root}, since it would have to have a fourth \gls{root} to form a \gls{conjugate_pair}.
            Thus, the third \gls{root} of a \gls{cubic} must a \gls{real_number}.
        \newpage

    \section{Argand Diagrams} \label{argand_diagrams}
        \subsection{Complex Numbers on Argand Diagrams}
            An \gls{argand_diagram} can be thought of similarly to a \gls{cartesian_graph}; \glspl{complex_number} in the form $x+iy$ can be plotted like $(x, y)$ can on a \gls{cartesian_graph}.
            Instead of an $x$ axis and a $y$ axis, there are instead a real axis and an imaginary axis, used to plot the \gls{real_part} and \gls{imaginary_part} of \glspl{complex_number}.
            An \gls{argand_diagram} looks like this:

            \begin{center}
            \begin{tikzpicture}[scale=1.5, >=stealth] % From Gemini
                % Draw axes
                \draw[->] (-0.5, 0) -- (4, 0) node[right] {$\real(z)$};
                \draw[->] (0, -0.5) -- (0, 3) node[above] {$\imag(z)$};
                
                % Define coordinates for complex number z = a + bi (e.g., z = 3 + 2i)
                \coordinate (O) at (0,0);
                \coordinate (Z) at (3,2);
                \coordinate (X) at (3,0);
                
                % Draw vector from origin to z
                %\draw[thick, red] (O) -- (Z) node[midway, above left] {$z = a+bi$};
                \draw[thick, red] (O) -- (Z);
                \draw[red] (Z) node[above] {$z=a+bi$};
                % Projection lines to axes
                \draw[dashed, gray] (Z) -- (X) node[below] {$a$};
                \draw[dashed, gray] (Z) -- (0,2) node[left] {$b$};
                
                % Plot point at z
                \filldraw[red] (Z) circle (1.5pt);
                
                % Draw angle theta
                %\draw (0.5,0) arc[start angle=0, end angle=33.7, radius=0.5];
                %\node at (0.7, 0.25) {$\theta$};
            \end{tikzpicture}
            \end{center}

            Where the point $z=a+bi$ is drawn where the point $(a, b)$ would be drawn on a \gls{cartesian_graph}.
            
        \subsection{Modulus-Argument Form}
            The \gls{modulus_argument_form} of a \gls{complex_number} is similar to the polar form of a point on a \gls{cartesian_graph}.
            It defines the point using its distance from the origin, and the angle it makes with the positive half of the real axis, in \glspl{radian}.
            The \gls{modulus} of $z$ is written as $|z|$, which is calculated using Pythagoras' Theorem:
            \begin{equation}
                |z|=\sqrt{\real^2(z)+\imag^2(z)}
            \end{equation}
            where $\real^2(z)$ and $\imag^2(z)$ mean $(\real(z))^2$ and $(\imag(z))^2$ respectively.
            The angle that $z$ makes with the positive real number line is called the \gls{argument}, and is positive, going anti-clockwise if $\imag(z)\ge0$, and negative, going clockwise if $\imag(z)<0$.
            The \gls{argument} of $z$ is written as $\arg(z)$.

            An \gls{argand_diagram} is split into four \glspl{quadrant}, called \gls{quadrant} 1, 2, 3 and 4:

            \begin{center}
            \begin{tikzpicture}[scale=1.5, >=stealth]
                \draw[->] (-2, 0) -- (2, 0) node[right] {$\real(z)$};
                \draw[->] (0, -2) -- (0, 2) node[above] {$\imag(z)$};
                
                \coordinate (Q1) at (1,1);
                \coordinate (Q2) at (-1,1);
                \coordinate (Q3) at (-1,-1);
                \coordinate (Q4) at (1,-1);
                
                \draw (Q1) node {Quadrant 1};
                \draw (Q2) node {Quadrant 2};
                \draw (Q3) node {Quadrant 3};
                \draw (Q4) node {Quadrant 4};
            \end{tikzpicture}
            \end{center}


            For values of $z$ in the first and fourth \glspl{quadrant} ($\real(z)>0$):
            \begin{equation}
                \arg(z)=\arctan\left(\frac{\imag(z)}{\real(z)}\right)
            \end{equation}
            For values of $z$ in the second \gls{quadrant}:
            \begin{equation}
                \arg(z)=\arctan\left(\frac{\imag(z)}{\real(z)}\right)+\pi
            \end{equation}
            For values of $z$ in the third \gls{quadrant}:
            \begin{equation}
                \arg(z)=\arctan\left(\frac{\imag(z)}{\real(z)}\right)-\pi
            \end{equation}

            On an \gls{argand_diagram}, the \gls{complex_conjugate} of a number is itself mirror in the line $\imag(z)=0$ (see subsection \ref{ssec:loci_and_regions_in_argand_diagrams}):

            \begin{center}
            \begin{tikzpicture}[scale=1.5, >=stealth]
                \draw[->] (-4, 0) -- (4, 0) node[right] {$\real(z)$};
                \draw[->] (0, -3) -- (0, 3) node[above] {$\imag(z)$};
                
                \coordinate (O) at (0,0);
                \coordinate (Z) at (3,2);
                \coordinate (conj_Z) at (3,-2);
                \coordinate (neg_Z) at (-3,-2);
                \coordinate (neg_conj_Z) at (-3,2);
                
                \draw[thick, red] (O) -- (Z);
                \draw[thick, red] (O) -- (conj_Z);
                \draw[thick, red] (O) -- (neg_Z);
                \draw[thick, red] (O) -- (neg_conj_Z);

                \draw[red] (Z) node[above] {$z=a+bi$};
                \draw[red] (conj_Z) node[below] {$z^*=a-bi$};
                \draw[red] (neg_Z) node[below] {$-z=-a-bi$};
                \draw[red] (neg_conj_Z) node[above] {$-z^*=-a+bi$};

                \filldraw[red] (Z) circle (1.5pt);
                \filldraw[red] (conj_Z) circle (1.5pt);
                \filldraw[red] (neg_Z) circle (1.5pt);
                \filldraw[red] (neg_conj_Z) circle (1.5pt);
            \end{tikzpicture}
            \end{center}

        \subsection{Loci and Regions in Argand Diagrams} \label{ssec:loci_and_regions_in_argand_diagrams}

        \newpage

    \section{Series I} \label{series_i}

    \section{Roots of Polynomials} \label{roots_of_polynomials}
        \subsection{Sums of Roots}
            The \glspl{root} of \glspl{polynomial} can be related using Vieta's formulae (see ).
            The \glspl{root} of a general \gls{polynomial} are written as $\alpha$, $\beta$, $\gamma$, $\delta$, etc., and the coefficients are written as $a$, $b$, $c$, $d$, $e$, etc., where $a$ is the coefficient of the greatest power of the \gls{polynomial}, $b$ is the coefficient of the next greatest power, etc.
            The sum of all \glspl{root} is then written as $\sum\alpha\equiv\alpha+\beta+\cdots$.
            Multiple \glspl{root} can be multiplied together and then summed, for example the sum of all pairs of \glspl{root} ($\alpha\beta+\alpha\gamma+\cdots$) would be written as $\sum\alpha\beta$.
            This would be the sum of \textit{all} possible pairs of \glspl{root}.

        \subsection{Linear Transformations of Roots}

    \newpage
    \newpage
    \section{Volumes of Revolution} \label{volumes_of_revolution}
    \newpage
    \section{Matrices} \label{matricies}
    \newpage
    \section{Linear Transformations} \label{linear_transformations}
    \newpage
    \section{Proof by Induction} \label{proof_by_induction}
    \newpage
    \section{Vectors} \label{vectors}
    \newpage

\chapter{Core Pure Year II} \label{ch:core_pure_ii}
    \section{Complex Numbers II} \label{complex_numbers_ii}
    \newpage
    \section{Series II} \label{series_ii}
    \newpage
    \section{Methods in Calculus} \label{methods_in_calculus}
    \newpage
    \section{Polar Coordinates} \label{polar_coordinates}
    \newpage
    \section{Hyperbolic Function} \label{hyperbolic_functions}
    \newpage
    \section{Methods in Differential Equations} \label{methods_in_differential_equations}
    \newpage
    \section{Modelling with Differential Equations} \label{nodelling_with_differential_equations}
    \newpage

\chapter{Further Mechanics Year I} \label{ch:mechanics_i}
    \section{Momentum and Impulse} \label{momentum_and_impulse}
    \newpage
    \section{Work, Energy and Power} \label{work_energy_and_power}
    \newpage
    \section{Elastic Strings and Springs} \label{elastic_strings_and_springs}
    \newpage
    \section{Elastic Collisions in One Dimension} \label{elastic_collisions_in_one_dimension}
    \newpage
    \section{Elastic Collisions in Two Dimensions} \label{elastic_collisions_in_two_dimensions}
    \newpage

\chapter{Further Mechanics Year II} \label{ch:mechanics_ii}
    \section{Circular Motion} \label{circular_motions}
    \newpage
    \section{Centre of Mass of Plane Figures} \label{centre_of_mass_of_plane_figures}
    \newpage
    \section{Further Centres of Mass} \label{further_centres_of_mass}
    \newpage
    \section{Kinematics} \label{kinematics}
    \newpage
    \section{Dynamics} \label{dynamics}
    \newpage

\chapter{Further Statistics Year I} \label{ch:statistics_i}
    \section{Discrete Random Variables} \label{}
    \newpage
    \section{Poisson Distributions} \label{}
    \newpage
    \section{Geometric and Negative Binomial Distributions} \label{}
    \newpage
    \section{Hypothesis Testing} \label{}
    \newpage
    \section{Central Limit Theorem} \label{}
    \newpage
    \section{$\chi^2$ Tests} \label{}
    \newpage
    \section{Probability Generating Functions} \label{}
    \newpage
    \section{Quality of Tests} \label{}
    \newpage

\chapter{Further Statistics Year II} \label{ch:statistics_ii}
    \section{Linear Regression} \label{}
    \newpage
    \section{Correlation} \label{}
    \newpage
    \section{Continuous Distribution} \label{}
    \newpage
    \section{Combinations of Random Variables} \label{}
    \newpage
    \section{Estimation, Confidence Intervals and Tests Using a Normal Distribution} \label{}
    \newpage
    \section{Further Hypothesis Tests} \label{}
    \newpage
    \section{Confidence Intervals and Tests Using the $t$-Distribution} \label{}
    \newpage

\chapter{Decision Year I} \label{ch:decision_i}
    \section{Algorithms} \label{algorithms}
        \section{Basic Algorithms}

        \newpage

        \section{Bubble Sort}
            \subsection{Description}
                Bubble sort is an algorithm that sorts an unsorted list of data.
                It passes through the list, checking each pair of items to see if they need to be swapped to be in order.
                For example, if it were sorting a list of numbers into numerical order and it found the pair 2, 1, it would swap them.
                This is an example of a single bubble sort pass, with each comparison marked:
                \begin{verbatim}
                    [3   5]  8   2   5   6   3
                    3  [5   8]  2   5   6   3
                    3   5  [8   2]  5   6   3
                    3   5   2  [8   5]  6   3
                    3   5   2   5  [8   6]  3
                    3   5   2   5   6  [8   3]
                    3   5   2   5   6   3   8
                \end{verbatim}
                Now, the 8 is in the right place, and the bubble sort will do another pass, ignoring the 8, since it must be in the right place after 1 pass, since it is the largest number.
                In general, after $n$ passes, the last $n$ items are sorted.
                Bubble sort gets its name from this property of the largets elements `bubbling' to the top (end) of the list.
                If the algorithm does a pass without swapping any items, it must have already finished sorting the list, so it will return early.

                Example of bubble sort: \\
                \begin{center}
                \begin{tabular}{c c c c c c c c}
                    Pass 1: &3 &5 &8 &2 &5 &6 &3 \\
                    Pass 2: &3 &5 &2 &5 &6 &3 &\textbf{8} \\
                    Pass 3: &3 &2 &5 &5 &3 &\textbf{6} &\textbf{8} \\
                    Pass 4: &2 &3 &5 &3 &\textbf{5} &\textbf{6} &\textbf{8} \\
                    Pass 6: &2 &3 &3 &\textbf{5} &\textbf{5} &\textbf{6} &\textbf{8} \\
                    Pass 7: &2 &3 &\textbf{3} &\textbf{5} &\textbf{5} &\textbf{6} &\textbf{8} \\
                    \multicolumn{8}{l}{\textbf{No more swaps; sort complete.}}
                \end{tabular}
                \end{center}
                Notice that pass 6 is the same as pass 7, as only on pass 7 has the algorithm not performed any swaps.

            \subsection{Order}
                \begin{theorem}
                    Bubble sort has an order of $\bigO(n)$
                \end{theorem}
                \begin{proof}
                    Let $n$ be the number of elements in the list to be sorted.
                    This means that (in the worst-case scenario) the bubble sort will have to do $i-1$ pass to guarantee that every item has been sorted (since after the $i$th pass, $i$ items have been sorted, $n-i$ items will be left to sort, the $i$th item must also be sorted).
                    In the $i$th pass, $i-1$ items will have been sorted, so the total number of comparisons will be:
                    $$\sum_{i=1}^n(n-i)=\frac{n(n-1)}{2}=\frac{1}{2}n^2+\frac{1}{2}n$$
                    Taking the fastest-growing term of $\frac{1}{2}n^2+\frac{1}{2}n$ (the $\frac{1}{2}n^2$ term) and discarding the constant coefficient ($\frac{1}{2}$), this gives $n^2$, thus the order of the bubble sort algorithm is $\bigO(n^2)$.
                \end{proof}
            
            \subsection{Implementation}
                In pseudocode, this is an example implementation of a bubble sort algorithm:
                \begin{verbatim}
// Sorts a list using bubble sort; returns sorted list
// Example:
//     Input:  [3, 2, 7, 3, 5]
//     Output: [2, 3, 3, 5, 7]
//     Explanation: The unsorted input is sorted so the the numbers are returned in order
FUNC bubble_sort(list):            // Takes an unsorted list as input
    VAR n = list.UPPER_BOUND

    FOR (i = 0 TO n):                     // Goes through each pass
        VAR hasSwapped = FALSE
        FOR (j = 0 TO n - i):             // Goes through each value in the list
            IF (list[i] > list[j + 1]):   // Checks if items need swapping
                hasSwapped = TRUE
                VAR temp = list[j + 1]    // Swaps items
                list[j + 1] = list[i]
                list[i] = temp
            END IF
        END FOR

        IF (NOT hasSwapped):   // If it doesn't find any items to swap, the list must be sorted
            OUTPUT("No more swaps; sort complete.")
            RETURN list        // Return the sorted list early, since there are no more swaps
        END IF
    END FOR

    RETURN list                // Returns sorted list
END FUNC
                \end{verbatim}
        \newpage

        \section{Quick Sort}
            \subsection{Description}

            \subsection{Order}

            \subsection{Implementation}

        \newpage

        \section{Bin-Packing Algorithms}

        \newpage
        
        \section{Order of Algorithms}

        \newpage

    \newpage
    \section{Graphs} \label{graphs}
    \newpage
    \section{Algorithms on Graphs} \label{algorithms_on_graphs}
    \newpage
    \section{Route Inspection} \label{route_inspection}
    \newpage
    \section{The Travelling Salesman Problem} \label{the_travelling_salesman_problem}
    \newpage
    \section{Linear Programming} \label{linea_programming}
    \newpage
    \section{The Simplex Algorithm} \label{the_simplex_algorithm}
    \newpage
    \section{Critical Path Analysis} \label{critical_path_analysis}
    \newpage

\chapter{Decision Year II} \label{ch:decision_ii}
    \section{Transportation Problems} \label{transportation_problems}
    \newpage
    \section{Allocation Problems} \label{allocation_problems}
    \newpage
    \section{Flows in Networks I} \label{flows_in_networks_i}
    \newpage
    \section{Flows in Networks II} \label{flows_in_networks_ii}
    \newpage
    \section{Dynamic Programming} \label{dynamic_programming}
    \newpage
    \section{Game Theory} \label{game_theory}
    \newpage
    \section{Recurrence Relations} \label{recurrence_relations}
    \newpage
    \section{Decision Analysis} \label{decision_analysis}
    \newpage


\chapter{Further Pure Mathematics Year I} \label{ch:pure_i}
    \section{Vectors} \label{}
    \newpage
    \section{Conic Sections I} \label{}
    \newpage
    \section{Conic Sections II} \label{}
    \newpage
    \section{Inequalities} \label{}
    \newpage
    \section{The $t$-Formulae} \label{}
    \newpage
    \section{Taylor Series} \label{}
    \newpage
    \section{Methods in Calculus} \label{}
    \newpage
    \section{Numerical Methods} \label{}
    \newpage
    \section{Reducible Differential Equations} \label{}
    \newpage

\chapter{Further Pure Mathematics Year II} \label{ch:pure_ii}
    \section{Number Theory} \label{}
    \newpage
    \section{Groups} \label{}
    \newpage
    \section{Complex Numbers} \label{}
    \newpage
    \section{Recurrence Relations} \label{}
    \newpage
    \section{Matrix Algebra} \label{}
    \newpage
    \section{Integration Techniques} \label{}
    \newpage

\appendix
\chapter{Appendix: Formulae} \label{ch:formulae}
    \section{Formula Sheet} \label{formula_sheet}
        \subsection{Pure Mathematics}
            \subsubsection{Summations}
            \begin{equation}
                \sum_{r=1}^nr^2=\frac{1}{6}n(n+1)(2n+1)
            \end{equation}
            \begin{equation}
                \sum_{r=1}^nr^3=\frac{1}{6}n^2(n+1)^2
            \end{equation}

            \subsubsection{Matrix Transformations}
            Anticlockwise rotation through $\theta$ about $O$:
            \begin{equation}
                \begin{bmatrix}
                    \cos\theta & -\sin\theta \\
                    \sin\theta & \cos\theta
                \end{bmatrix}
            \end{equation}
            Reflection in the line $y=(\tan\theta)x$:
            \begin{equation}
                \begin{bmatrix}
                    \cos2\theta & \sin2\theta \\
                    \sin2\theta & -\cos2\theta
                \end{bmatrix}
            \end{equation}



            The matrices for rotations (in three dimensions) though an angle $\theta$ about one of the axes are:
            \begin{equation}
                \begin{bmatrix}
                    1 & 0 & 0 \\
                    0 & \cos\theta & -\sin\theta \\
                    0 & \sin\theta & \cos\theta
                \end{bmatrix}
                \quad\text{for the}\ x\text{-axis}
            \end{equation}
            \begin{equation}
                \begin{bmatrix}
                    \cos\theta & 0 & \sin\theta \\
                    0 & 1 & 0 \\
                    -\sin\theta & 0 & \cos\theta
                \end{bmatrix}
                \quad\text{for the}\ y\text{-axis}
            \end{equation}
            \begin{equation}
                \begin{bmatrix}
                    \cos\theta & -\sin\theta & 0 \\
                    \sin\theta & \cos\theta & 0 \\
                    0 & 0 & 1
                \end{bmatrix}
                \quad\text{for the}\ z\text{-axis}
            \end{equation}



            \subsubsection{Area of a Sector}
            \begin{equation}
                A=\frac{1}{2}\int r^2d\theta\quad\text{(polar coordinates)}
            \end{equation}

            \subsubsection{Complex Numbers}
            \begin{equation}
                [r\cos\theta+i\sin\theta]^n=r^n(\cos n\theta+i\sin n\theta)
            \end{equation}
            The \glspl{root} of $z^n=1$ are given by
            \begin{equation}
                z=e^{\frac{2\pi ki}{n}},\ \text{for}\ k=0,1,2,\dots,n-1
            \end{equation}

            \subsubsection{Maclaurin's and Taylor's Series}
            \begin{equation}
                f(x)=f(0)+xf'(0)+\frac{x^2}{2!}f''(0)+\cdots+\frac{x^r}{r!}f^{(r)}(0)+\cdots
            \end{equation}
            \begin{equation}
                e^x=\exp(x)=1+x+\frac{x^2}{2!}+\cdots+\frac{x^r}{r!}+\cdots\quad\text{for all}\ x
            \end{equation}
            \begin{equation}
                \ln(n+1)=x-\frac{x^2}{2}+\frac{x^3}{3}-\cdots+(-1)^r\frac{x^r}{r}+\cdots\quad(-1<x\le1)
            \end{equation}
            \begin{equation}
                \cos x=1-\frac{x^3}{3!}+\frac{x^5}{5!}-\cdots+(-1)^{r+1}\frac{x^{2r+1}}{(2r+1)!}+\cdots\quad\text{for all}\ x
            \end{equation}
            \begin{equation}
                \cos x=1-\frac{x^2}{2!}+\frac{x^4}{4!}-\cdots+(-1)^r\frac{x^{2r}}{(2r)!}\quad\text{for all}\ x
            \end{equation}
            \begin{equation}
                \arctan x=x-\frac{x^3}{3}+\frac{x^5}{5}-\cdots+(-1)^r\frac{x^{2r+1}}{2r+1}+\cdots\quad(-1\le x\le1)
            \end{equation}

            \subsubsection{Vectors}
            The resolved part of $\textbf{a}$ in the direction of $\textbf{b}$ is
            \begin{equation}
                \frac{\textbf{a}\cdot\textbf{b}}{\textbf{b}}
            \end{equation}
            Vector Product
            \begin{equation}
                \textbf{a}\times\textbf{b}=|\textbf{a}||\textbf{b}|\sin\theta\hat{\textbf{n}}=\left|
                    \begin{matrix}
                        \textbf{i} & \textbf{j} & \textbf{k} \\
                        a_1 & a_2 & a_3 \\
                        b_1 & b_2 & b_3 \\
                    \end{matrix}
                \right|=
                \begin{bmatrix}
                    a_2b_3-a_3b_2 \\
                    a_3b_1-a_1b_3 \\
                    a_1b_2-a_2b_1
                \end{bmatrix}
            \end{equation}
            \begin{equation}
                \textbf{a}\cdot(\textbf{b}b\times\textbf{c})=\left|
                    \begin{matrix}
                        a_1 & a_2 & a_3 \\
                        b_1 & b_2 & b_3 \\
                        c_1 & c_2 & c_3
                    \end{matrix}
                \right|=\textbf{b}(\textbf{c}\times\textbf{a})=\textbf{c}=(\textbf{a}\times\textbf{b})
            \end{equation}
            If $A$ is the point with position vector $\textbf{a}=a_1\textbf{i}+a_2\textbf{j}+a_3\textbf{k}$ and the direction vector $\textbf{b}$ is given by $\textbf{b}=b_1\textbf{i}+b_2\textbf{j}+b_3\textbf{k}$, then the straight line through $A$ with direction vector $\textbf{b}$ has Cartesian equation:
            \begin{equation}
                \frac{x-a_1}{b_1}=\frac{y-a_2}{b_2}=\frac{z-a_3}{b_3}=\lambda
            \end{equation}
            The plane through $A$ with normal vector $\textbf{n}=n_1\textbf{i}+n_2\textbf{j}+n_3\textbf{k}$ has Cartesian equation:
            \begin{equation}
                n_1x+n_2y+n_3z+d=0\text{ where }d=-\textbf{a}\cdot\textbf{n}
            \end{equation}
            The plane through non-collinear points $A$, $B$ and $C$ has vector equation:
            \begin{equation}
                \textbf{r}=\textbf{a}+\lambda(\textbf{b}-\textbf{a})+\mu(\textbf{c}-\textbf{a})=(1-\lambda-\mu)\textbf{a}+\lambda\textbf{b}+\mu\textbf{c}
            \end{equation}
            The plane through $A$ and parallel to $\textbf{b}$ and $\textbf{c}$ as vector equation:
            \begin{equation}
                \textbf{r}=\textbf{a}+s\textbf{b}+t\textbf{c}
            \end{equation}
            The perpendicular distance of $(\alpha,\beta,\gamma)$ from $n_1x+n_2y+n_3z+d=0$ is:
            \begin{equation}
                \frac{|n_1\alpha+n_2\beta+n_3\gamma+d|}{\sqrt{n_1^2+n_2^2+n_3^2}}
            \end{equation}

            \subsubsection{Hyperbolic Functions}
            \begin{equation}
                \cosh^2x-\sinh^2x=1
            \end{equation}
            \begin{equation}
                \sinh2x=2\sinh x\cosh x
            \end{equation}
            \begin{equation}
                \cosh2x=\cosh^2x+\sinh^2x
            \end{equation}
            \begin{equation}
                \cosh^{-1}=\ln\left\{x+\sqrt{x^2-1}\right\}\quad{(x\ge1)}
            \end{equation}
            \begin{equation}
                \sinh^{-1}=\frac{1}{2}\ln\left(\frac{1+x}{1-x}\right)\quad{(|x|<1)}
            \end{equation}

            \subsubsection{Conics}
            \noindent\textbf{Ellipses} \\

            Standard Form:
            \begin{equation}
                \frac{x^2}{a^2}+\frac{y^2}{b^2}=1
            \end{equation}
            Parametric Form:
            \begin{equation}
                (a\cos\theta,b\sin\theta)
            \end{equation}
            Eccentricity
            \begin{equation}
                \begin{gathered}
                    e<1 \\ b^2=a^2(1-e^2)
                \end{gathered}
            \end{equation}
            Foci:
            \begin{equation}
                (\pm ae,0)
            \end{equation}
            Directrices:
            \begin{equation}
                x=\pm\frac{a}{e}
            \end{equation}
            Asymptotes: none \\

            \noindent\textbf{Parabolas} \\

            Standard Form:
            \begin{equation}
                y=4ax
            \end{equation}
            Parametric Form:
            \begin{equation}
                (at^2,2at)
            \end{equation}
            Eccentricity:
            \begin{equation}
                e=1
            \end{equation}
            Foci:
            \begin{equation}
                (a,0)
            \end{equation}
            Directrices:
            \begin{equation}
                x=-a
            \end{equation}
            Asymptotes: none \\

            \noindent\textbf{Hyperbolae} \\

            Standard Form:
            \begin{equation}
                \frac{x^2}{a^2}-\frac{y^2}{b^2}=1
            \end{equation}
            Parametric Forms:
            \begin{equation}
                (a\sec\theta,b\tan\theta)
            \end{equation}
            \begin{equation}
                (\pm a\cosh\theta,b\sinh\theta)
            \end{equation}
            Eccentricity:
            \begin{equation}
                \begin{gathered}
                    e>1 \\
                    b^2=a^2(e^2-1)
                \end{gathered}
            \end{equation}
            Foci:
            \begin{equation}
                (\pm ae,0)
            \end{equation}
            Directrices:
            \begin{equation}
                x=\pm\frac{a}{e}
            \end{equation}
            Asymptotes:
            \begin{equation}
                \frac{x}{a}=\pm\frac{y}{b}
            \end{equation}

            \noindent\textbf{Rectangular Hyperbolae} \\

            Standard Form:
            \begin{equation}
                xy=c^2
            \end{equation}
            Parametric Form:
            \begin{equation}
                \left(ct,\frac{c}{t}\right)
            \end{equation}
            Eccentricity:
            \begin{equation}
                e=\sqrt{2}
            \end{equation}
            Foci:
            \begin{equation}
                (\pm\sqrt{2}c,\pm\sqrt{2}c)
            \end{equation}
            Directrices:
            \begin{equation}
                x+y=\pm\sqrt{2}c
            \end{equation}
            Asymptotes:
            \begin{equation}
                x=0,\ y=0
            \end{equation}

            \subsubsection{Differentiation}
            \begin{alignat}{2}
                f(x)&\qquad f'(x) \notag\\
                \arcsin x&\qquad\frac{1}{\sqrt{1-x^2}} \\
                \arccos x&\qquad-\frac{1}{\sqrt{1-x^2}} \\
                \arctan x&\qquad\frac{1}{\sqrt{1+x^2}} \\
                \sinh x&\qquad\cosh x \\
                \cosh x&\qquad\sinh x \\
                \tanh x&\qquad\sech^2x \\
                \arsinh x&\qquad\frac{1}{1+x^2} \\
                \arcosh x&\qquad\frac{1}{\sqrt{x^2-1}} \\
                \artanh x&\qquad \frac{1}{1-x^2}\\
            \end{alignat}

            \subsubsection{Integration (+ constant; $a>0$ where relevant)}
            \begin{alignat}{2}
                f(x)&\qquad \int f(x)dx \notag\\
                \sinh x&\qquad\cosh x \\
                \cosh x&\qquad\sinh x \\
                \tanh x&\qquad\ln\cosh x \\
                \frac{1}{\sqrt{a^2-x^2}}&\qquad\sin^-1\left(\frac{x}{a}\right)+c\quad(|x|<a) \\
                \frac{1}{a^2+x^2}&\qquad\frac{1}{a}\tan^-1\left(\frac{x}{a}\right)+c \\
                \frac{1}{x^2-a^2}&\qquad\cosh^-1(\frac{x}{a})\ \text{or}\ \ln\left\{x+\sqrt{x^2-a^2}+c\quad(x>a)\right\} \\
                \frac{1}{\sqrt{a^2+x^2}}&\qquad\sinh^{-1}\left(\frac{x}{a}\right)\ \text{or}\ \ln\left\{x+\sqrt{x^2+a^2}\right\}+c \\
                \frac{1}{a^2-x^2}&\qquad\frac{1}{2a}\ln\left|\frac{a+x}{a-x}\right|=\frac{1}{a}\tanh^{-1}\left(\frac{x}{a}\right)+c\quad(|x|<a) \\
                \frac{1}{x^2-a^2}&\qquad\frac{1}{2a}\ln\left|\frac{a-x}{a+x}\right|+c
            \end{alignat}

            \subsubsection{Arc Length}
            \begin{equation}
                s=\int\sqrt{1+\left(\frac{dy}{dx}\right)^2}dx\quad\text{(Cartesian coordinates)}
            \end{equation}
            \begin{equation}
                s=\int\sqrt{\left(\frac{dx}{dt}\right)^2+\left(\frac{dy}{dt}\right)^2}dt\quad\text{(parametric form)}
            \end{equation}
            \begin{equation}
                s=\int\sqrt{r^2+\left(\frac{dr}{d\theta}\right)^2}d\theta\quad\text{(polar form)}
            \end{equation}

            \subsubsection{Surface Area of Revolution}
            \begin{equation}
                s_x=2\pi\int y\sqrt{1+\left(\frac{dy}{dx}\right)^2}dx\quad\text{(Cartesian coordinates)}
            \end{equation}
            \begin{equation}
                s_x=2\pi\int y\sqrt{\left(\frac{dx}{dt}\right)^2+\left(\frac{dy}{dt}\right)^2}dt\quad\text{(parametric form)}
            \end{equation}
            \begin{equation}
                s_x=2\pi\int r\sin\theta\sqrt{r^2+\left(\frac{dr}{d\theta}\right)^2}d\theta\quad\text{(polar form)}
            \end{equation}

            \subsection{Mechanics}
            \subsubsection{Constant Acceleration}
            \begin{alignat}{2}
                s=ut+\frac{1}{2}at^2&\qquad\textbf{s}=\textbf{u}t+\frac{1}{2}\textbf{a}t^2 \\
                s=vt-\frac{1}{2}at^2&\qquad\textbf{s}=\textbf{v}t-\frac{1}{2}\textbf{a}t^2 \\
                v=u+at&\qquad\textbf{v}=\textbf{u}+\textbf{a}t \\
                s=\frac{1}{2}(u+v)t&\qquad\textbf{s}=\frac{1}{2}(\textbf{u}+\textbf{v})t \\
                v^2=u^2+2as&
            \end{alignat}

            \subsubsection{Centres of Mass}
            For uniform bodies: \\
            \begin{equation}
                \text{Triangular lamina:}\ \frac{2}{3}\ \text{along median from vertex}
            \end{equation}
            \begin{equation}
                \text{Solid hemisphere, radius }r\text{:}\ \frac{3}{8}r\ \text{from centre}
            \end{equation}
            \begin{equation}
                \text{Hemispherical shell:}\ \frac{2}{3}\ \text{along median from vertex}
            \end{equation}
            \begin{equation}
                \text{Circular arc, radius }r\text{ angle at centre }2\alpha\text{:}\ \frac{r\sin\alpha}{\alpha}\ \text{from centre}
            \end{equation}
            \begin{equation}
                \text{Sector of circle, radius }r\text{ angle at centre }2\alpha\text{:}\ \frac{2r\sin\alpha}{3\alpha}\ \text{from centre}
            \end{equation}
            \begin{equation}
                \text{Solid cone or pyramid of height }h\text{:}\ \frac{1}{4}h\ \text{above the base on the line from the centre of the base to the vertex}
            \end{equation}
            \begin{equation}
                \text{Conical shell of height }h\text{:}\ \frac{1}{3}h\ \text{above the base on the line from the centre of the base to the vertex}
            \end{equation}
        \newpage

        \subsection{Probability and Statistics}
        \subsubsection{Probability}
        \begin{equation}
            P(A\cup B)=P(A)+P(B)-P(A\cap B)
        \end{equation}
        \begin{equation}
            P(A\cap B)=P(A)\times P(B|A)
        \end{equation}

        \subsubsection{Standard Deviation}
        \begin{equation}
            \sqrt{\frac{\sum(x-\bar{x})^2}{n}}=\sqrt{\frac{\sum x^2}{n}-\bar{x}^2}
        \end{equation}

        \subsubsection{Discrete Distributions}
        \noindent\textbf{Binomial $\binomial(n,p)$} \\
        $P(X=x)$:
        \begin{equation}
            \binom{n}{x}p^x(1-p)^{n-x}
        \end{equation}
        Mean:
        \begin{equation}
            np
        \end{equation}
        Variance:
        \begin{equation}
            np(1-p)
        \end{equation}

        \noindent\textbf{Poisson $\poisson(\lambda)$} \\
        $P(X=x):$
        \begin{equation}
            e^{-\lambda}\frac{\lambda^x}{x!}
        \end{equation}
        Mean:
        \begin{equation}
            \lambda
        \end{equation}
        Variance:
        \begin{equation}
            \lambda
        \end{equation}
        
        \paragraph{Sampling Distributions}
        For a random sample $X_1,X_2,\dots,X_n$ of $n$ independent observations from a distribution having mean $\mu$ and a variance $\sigma^2$:
        \begin{equation}
            \bar{X}\text{ is an unbiased estimator of }\mu\text{, with }\Var(\bar{X})=\frac{\sigma^2}{n}
        \end{equation}
        \begin{equation}
            S^2\text{ is an unbiased estimator of }\sigma^2\text{, where }S^2=\frac{\sum(X_i-\bar{X})^2}{n-1}
        \end{equation}
        For a random sample of $n$ observations from $\normal(\mu,\sigma^2)$:
        \begin{equation}
            \frac{\bar{X}-\mu}{\sigma}\sim\normal(0,1)
        \end{equation}
        \begin{equation}
            \frac{\bar{X}-\mu}{\frac{S}{\sqrt{n}}}\sim t_{n-1}
        \end{equation}

        \subsubsection{Distribution-Free (Non-Parametric) Tests}
        Contingency tables:
        \begin{equation}
            \sum \frac{(O_i-E_i)^2}{E_i}\text{ is approximately distributed as }\chi^2
        \end{equation}
        \newpage

    \section{Non-Formula Sheet} \label{non_formula_sheet}
        \subsection{Pure Mathematics}
            Quadratic Formula
            \begin{equation} \label{eq:quadratic_formula}
                x=\frac{-b\pm\sqrt{b^2-4ac}}{2a}
            \end{equation}
        
            \subsubsection{Vieta's Formulae}
                For a \gls{quadratic} $ax^2+bx+c$ with \glspl{root} $\alpha$, $\beta$:
                \begin{equation}
                    \sum\alpha=\alpha+\beta=-\frac{b}{a}
                \end{equation}
                \begin{equation}
                    \alpha\beta=\frac{c}{a}
                \end{equation}

                For a cubic $ax^3+bx^2+cx+d$ with \glspl{root} $\alpha$, $\beta$, $\gamma$:
                \begin{equation}
                    \sum\alpha=\alpha+\beta+\gamma=-\frac{b}{a}
                \end{equation}
                \begin{equation}
                    \sum\alpha\beta=\alpha\beta+\alpha\gamma+\beta\gamma=\frac{c}{a}
                \end{equation}
                \begin{equation}
                    \alpha\beta\gamma=-\frac{d}{a}
                \end{equation}

                For a quartic $ax^4+bx^3+cx^2+dx+e$ with \glspl{root} $\alpha$, $\beta$, $\gamma$, $\delta$:
                \begin{equation}
                    \sum\alpha=\alpha+\beta+\gamma+\delta=-\frac{b}{a}
                \end{equation}
                \begin{equation}
                    \sum\alpha\beta=\alpha\beta+\alpha\gamma+\alpha\delta+\beta\gamma+\beta\delta+\gamma\delta=\frac{c}{a}
                \end{equation}
                \begin{equation}
                    \sum\alpha\beta\gamma=\alpha\beta\gamma+\alpha\beta\delta+\alpha\gamma\delta+\beta\gamma\delta=-\frac{d}{a}
                \end{equation}
                \begin{equation}
                    \alpha\beta\gamma\delta=\frac{e}{a}
                \end{equation}


                \subsubsection{{Integration}}
                \begin{equation}
                    \int u\frac{du}{dx}dx=uv-\int v\frac{du}{dx}dx
                \end{equation}
                \begin{equation}
                    \int \frac{f'(x)}{f(x)}=\ln|f(x)|+c
                \end{equation}
                \begin{align}
                    \tan x&\qquad\ln|\sec x|+c \\
                    \cot x&\qquad\ln|\sin x|+c \\
                    \cosec x&\qquad-\ln|\cosec x+\cot x|=\ln\left|\tan \frac{1}{2}x\right|+c \\
                    \sec x&\qquad\ln|\sec x+\tan x|=\ln\left|\tan\left(\frac{1}{2}x+\frac{1}{4}\pi\right)\right|+c \\
                    \sec^2x&\qquad\tan x+c \\
                    \tanh x&\qquad\ln\cosh x+c
                \end{align}

            \subsubsection{Binomial Series}
                \begin{equation}
                    (a+b)^n=a^n+\binom{n}{1}a^{n-1}b+\binom{n}{2}a^{n-2}b^2+\cdots+\binom{n}{r}a^{n-r}b^r+\cdots+b^n\quad(n\in\N)
                \end{equation}

                \subsubsection{Arithmetic Series}
                \begin{equation} \label{eq:arithmetic_series_finite_sum}
                    S_n=\frac{1}{2}n(a+l)=\frac{1}{2}n[2a+(n-1)d]
                \end{equation}

                \subsubsection{Geometric Series}
                \begin{equation} \label{eq:geometric_series_finite_sum}
                    S_n=\frac{a(1-r^n)}{1-r}
                \end{equation}
                \begin{equation} \label{eq:geometric_series_infinite_sum}
                    S_\infty=\frac{a}{1-r}
                \end{equation}

            \subsubsection{Trigonometry: Small Angles}
                \begin{equation} \label{sine_small_angle}
                    \sin\theta\approx\theta
                \end{equation}
                \begin{equation} \label{cosine_small_angle}
                    \cos\theta\approx1-\frac{\theta^2}{2}
                \end{equation}
                \begin{equation} \label{tangent_small_angle}
                    \tan\theta\approx\theta
                \end{equation}

            \subsubsection{Trigonometric Identities}
                \begin{equation} \label{sine_compound_angle_formula}
                    \sin(A\pm B)=\sin A\cos B\pm\cos A\sin B
                \end{equation}
                \begin{equation} \label{cosine_compound_angle_formula}
                    \cos(A\pm B)=\cos A\cos B\mp\sin A\sin B
                \end{equation}
                \begin{equation} \label{sine_addition_formula}
                    \sin A+\sin B=2\sin{\frac{A+B}{2}}\cos{\frac{A-B}{2}}
                \end{equation}
                \begin{equation} \label{sine_subtraction_formula}
                    \sin A-\sin B=2\cos{\frac{A+B}{2}}\sin{\frac{A-B}{2}}
                \end{equation}
                \begin{equation} \label{cosine_addition_formula}
                    \cos A+\cos B=2\cos{\frac{A+B}{2}}\cos{\frac{A-B}{2}}
                \end{equation}
                \begin{equation} \label{cosine_subtraction_formula}
                    \cos A-\cos B=-2\sin{\frac{A+B}{2}}\sin{\frac{A-B}{2}}
                \end{equation}

            \subsubsection{Differentiation from First Principles}
                \begin{equation}
                    f'(x)=\lim{h->0}{\frac{f(x+h)-f(x)}{h}}
                \end{equation}

            \subsubsection{Numerical Solutions of Equations}
                The Newton-Raphson iteration for solving $f(x)=0$:
                \begin{equation}
                    x_(n+1)=x_n-\frac{f(x_n)}{f'(x_n)}
                \end{equation}

            \subsubsection{Numerical Integration}
                The Trapezium Rule:
                \begin{equation}
                    \int_a^bydx\approx\frac{1}{2}\left\{h(y_0+y_n)+2(y_1+y_2+\cdots+y_{n-1})\right\},\ \text{where}\ h=\frac{b-a}{n}
                \end{equation}
            
            \subsubsection{Further Numerical Integration}
                The mid-ordinate rule:
                \begin{equation}
                    \int_a^bydx\approx h\left(y_{\frac{1}{2}+y_{\frac{3}{2}+\cdots+y_{n-\frac{3}{2}}+y_{n-\frac{1}{2}}}}\right)
                \end{equation}
                where $h=\frac{b-a}{n}$ \\
                Simpson's Rule:
                \begin{equation}
                    \int_a^bydx\approx\frac{1}{3}h\left\{(y_0+y_n)+4(y_1+y_3+\cdots+y_{n-1}+2(y_2+y_4+\cdots+y_{n-2}))\right\}
                \end{equation}

            \subsubsection{Numerical Solutions to Differential Equations}
                For $\frac{dy}{dx}=f(x)$ and small $h$, recurrence relations are: \\
                Euler's Method:
                \begin{equation}
                    y_{n+1}=y_n+hf(x),\quad x_{n+1}=x_n+h
                \end{equation}
                For $\frac{dy}{dx}=f(x,y)$: \\
                Euler's Method:
                \begin{equation}
                    y_{r+1}=y_r+hf(x_r,y_r),\quad x_{r+1}=x^r+h
                \end{equation}
                Improved Euler's Method:
                \begin{equation}
                    y_{r+1}=y_{r-1}+2hf(x_r,y_r),\quad x_{r+1}=x_r+h
                \end{equation}

            \subsection{Mechanics}

        \subsection{Probability and Statistics}

    \newpage

\printnoidxglossary

\end{document}